\section{Лево(право)линейные грамматики}

Наложив некоторые ограничения на внешний вид правил грамматики можно получить грамматики, задающие регулярные языки.

\begin{definition}[Леволинейная грамматика]
    Грамматика $G=\langle \Sigma, N, P, S \rangle$ называется леволинейной, если все её правила имеют вид
    \[N_i \to \alpha w,\]
    где $N_i \in N$, $\alpha \in \{\varepsilon\} \cup N$, $w \in \Sigma ^*$.
\end{definition}

\begin{definition}[Праволинейная грамматика]
    Грамматика $G=\langle \Sigma, N, P, S \rangle$ называется праволинейной, если все её правила имеют вид
    \[N_i \to  w \alpha,\]
    где $N_i \in N$, $\alpha \in \{\varepsilon\} \cup N$, $w \in \Sigma ^*$.
\end{definition}

Ноам Хомский и Джордж Миллер в работе~\sidecite{chomsky1958finite} показали, что лево(право)линейные грамматики задают регулярные языки.
Приведём процедуры построения автомата по грамматике и наоборот, грамматики по автомату.

Пусть дан конечный автомат $M = \langle \Sigma, Q, q_s, Q_f, \delta \rangle$. По нему можно построить праволинейную грамматику $G=\langle \Sigma, N, S, P \rangle$, где
\begin{itemize}
    \item $N = Q$
    \item $P = \{ q_i \to t q_j \mid (q_i, t, q_j)\in \delta\} \cup \{ q_i \to \varepsilon \mid q_i \in Q_F\}$
    \item $S = q_s$
\end{itemize}

Аналогичным образом строится автомат по праволинейной грамматике.
Упростить процедуру можно если заранее привести правила к виду $N_i \to tN_j$, где $t\in \Sigma$, добавив необходимое количество новых нетерминалов:
правило вида $N_i \to twN_k$ преобразуется в два правила
\begin{align*}
    N_i & \to tN_l  \\
    N_l & \to wN_k,
\end{align*}
а если цепочка $w$ содержит более одного символа, то преобразование продолжается аналогичным образом для нетерминала $N_l$ (и далее), пока все правила не примут нужный вид.

\begin{example}[Построение грамматики по автомату]
Построим грамматику по автомату, представленному на рисунке~\ref{fig:dfa_example}.
В автомате три состояния, соответственно, в грамматике будет три нетерминала: $Q_0, Q_1, Q_2$.
Нетерминал $Q_0$ будет стартовым, так как он соответствует стартовому состоянию.
Правила грамматики формируются на основе переходов в автомате и правил для финальных состояний, и в нашем случае будут следующими:
\begin{align*}
    Q_0 \to a Q_1  \quad & \quad Q_1 \to a Q_1 \\
    Q_1 \to b Q_2  \quad & \quad Q_2 \to b Q_1 \\
                     \quad & \quad Q_1 \to \varepsilon.
\end{align*}
Таким образом, мы получили грамматику, которая порождает тот же язык, что и исходный автомат%
\sidenote{В качестве упражнения можно взять несколько цепочек, принадлежащих языку, и проследить, как переходы в автомате в процессе раскознования цепочки соотносятся с применением правил грамматики при её выоде.}.
\end{example}

\begin{marginfigure}
    \begin{center}
                \begin{tikzpicture}
            \node[state,initial] (S) {$S$};
            \node[state,accepting,right=of S] (N1) {$N_1$};
            \node[state,above=of N1] (N2) {$N_2$};

            \path[->]
                (S) edge[bend left, above]   node{$a$} (N1)
                (N1) edge[loop below, below] node{$a$} (N1)
                (N1) edge[bend left, left]   node{$b$} (N2)
                (N2) edge[bend left, right]  node{$b$} (N1);
        \end{tikzpicture}

    \end{center}
    \caption{Автомат, построенный по грамматике $S \to aN_1$, $N_1 \to aN_1 \mid bbN_1 \mid \varepsilon$.
    Эквивалентен автомату из примера~\ref{exmpl:DFA}.}
    \label{fig:grammar_to_dfa}
\end{marginfigure}

\begin{example}[Построение автомата по грамматике]
Рассмотрим грамматику:
\begin{align*}
    S \to a N_1, \quad
    N_1 \to a N_1 \mid b b N_1 \mid \varepsilon.
\end{align*}
Построим по ней конечный автомат.

Первым шагом разобьём продукцию $N_1 \to b b N_1$ на две, введя дополнительный нетерминал.
Новая грамматика будет иметь следующий вид.

\begin{align*}
    S \to a N_1, \quad
    N_1 \to a N_1 \mid b N_2 \mid \varepsilon , \quad
    N_2 \to b N_1.
\end{align*}

Теперь можно заняться конструированием автомата.
Множество состояний автомата соответствует нетерминалам грамматики:
\begin{align*}
    Q = \{S, N_1, N_2\}.
\end{align*}
Алфавит: $\Sigma = \{a, b\}$.
Переходы автомата определяются на основе правил грамматики:
\begin{align*}
\delta = \{
(S, a, N_1),\;
(N_1, a, N_1),\;
(N_1, b, N_2),\;
(N_2, b, N_1)
\}.
\end{align*}
Финальные состояния~--- те, для которых в грамматике есть правило с правой частью $\varepsilon$:
\begin{align*}
Q_F = \{ N_1 \}.
\end{align*}
Таким образом, получаем автомат
\begin{align*}
M = \langle \Sigma, Q, S, Q_F, \delta \rangle,
\end{align*}
распознающий язык всех цепочек, начинающихся с символа $a$ и содержащих произвольное количество $a$ и пар $bb$ после него.

Построенный автомат показан на рисунке~\ref{fig:grammar_to_dfa}. Обратите внимание, что этот автомат эквивалентен автомату из примера~\ref{exmpl:DFA}, только обозначения состояний изменены: $S \leftrightarrow 0$, $N_1 \leftrightarrow 1$, $N_2 \leftrightarrow 2$.
\end{example}
